\documentclass[11pt,a4paper]{article}

\usepackage[utf8]{inputenc}
\usepackage[T1]{fontenc}
\usepackage{amsmath,amssymb,amsfonts}
\usepackage{booktabs,tabularx,longtable}
\usepackage{enumitem}
\usepackage{geometry}
\usepackage{hyperref}
\usepackage{xcolor}
\usepackage{listings}
\usepackage{fancyhdr}
\usepackage{titlesec}
\usepackage{colortbl}
\usepackage{array}

\geometry{margin=2.2cm,top=2.8cm,bottom=2.8cm}

\definecolor{keepgreen}{HTML}{2E7D32}
\definecolor{removered}{HTML}{C62828}
\definecolor{modifyamber}{HTML}{EF6C00}
\definecolor{structblue}{HTML}{1565C0}
\definecolor{llmorange}{HTML}{E65100}
\definecolor{codebg}{HTML}{F5F5F5}

\hypersetup{
  colorlinks=true,
  linkcolor=blue!60!black,
  citecolor=blue!60!black,
  urlcolor=blue!60!black,
  pdftitle={ISLS Final Architecture --- Clean Implementation Spec},
  pdfauthor={Sebastian Klemm},
}

\lstset{
  basicstyle=\ttfamily\scriptsize,
  breaklines=true,
  frame=single,
  backgroundcolor=\color{codebg},
  keywordstyle=\color{blue!70!black}\bfseries,
  stringstyle=\color{red!60!black},
  commentstyle=\color{green!50!black}\itshape,
  numbers=left,
  numberstyle=\tiny\color{gray},
  numbersep=5pt,
  tabsize=4,
  showstringspaces=false,
  xleftmargin=1.5em,
  framexleftmargin=1em,
}

\titleformat{\section}{\Large\bfseries}{\thesection}{1em}{}
\titleformat{\subsection}{\large\bfseries}{\thesubsection}{1em}{}
\titleformat{\subsubsection}{\normalsize\bfseries}{\thesubsubsection}{1em}{}

\pagestyle{fancy}
\fancyhf{}
\fancyhead[L]{\small ISLS Final Architecture Spec}
\fancyhead[R]{\small Klemm, 2026}
\fancyfoot[C]{\thepage}

\begin{document}

% ============================================================
% TITLE PAGE
% ============================================================
\begin{titlepage}
\begin{center}
\vspace*{1.5cm}
{\Huge\bfseries ISLS}\\[0.4cm]
{\Huge\bfseries Final Architecture}\\[0.3cm]
{\LARGE\bfseries Clean Implementation Specification}\\[1.2cm]
{\Large D1: Pipeline $\to$ Compile $\to$ Docker $\to$ Browser}\\[1.5cm]
{\large Sebastian Klemm}\\[0.3cm]
{\small March 2026}\\[2cm]

\begin{tabular}{ll}
\textbf{Repository:} & \texttt{graphity/isls} \\
\textbf{Builds on:} & ISLS v3.4 (PCR-Conformant HDAG Codegen) \\
\textbf{Core engine:} & \texttt{isls-forge-llm} (Staged Closure) \\
\textbf{Default LLM:} & \texttt{gpt-4o} \\
\textbf{License:} & MIT \\
\end{tabular}

\vfill
{\small
\textbf{Purpose:} ONE spec. ONE implementation pass. No incremental patches.\\[0.2cm]
This document describes the complete target system, instructs Claude Code\\
to clean the workspace of all obsolete crates, and codifies the 12\\
architectural rules distilled from 8 debug cycles.\\[0.2cm]
After implementation: \texttt{cargo build} $\to$ \texttt{forge-v2} $\to$ \texttt{docker-compose up} $\to$ browser.\\[0.2cm]
\textbf{Clean.} No Blinddarm. No dead code. No template path. Only HDAG.
}
\end{center}
\end{titlepage}

\tableofcontents
\newpage

% ============================================================
\section{Situation}
% ============================================================

The ISLS workspace has accumulated ${\sim}25$ crates across versions v1.0--v3.4.
Two parallel code generation paths exist:

\begin{enumerate}
\item \textbf{Template path} (v2.x): \texttt{isls-orchestrator} with Tera templates + \texttt{isls-renderloop} multi-pass LLM enrichment.
\item \textbf{HDAG path} (v3.4): \texttt{isls-forge-llm} with PCR Staged Closure + deterministic structural generators + typed HDAG dependency graph.
\end{enumerate}

The HDAG path is architecturally correct. It has been validated through 8 iterative
cycles that systematically expanded the structural (deterministic) layer from 18 to 40
files and reduced LLM-generated files from 37 to 15. The template path is obsolete.
Both coexist in the workspace, causing confusion, compile-time waste, and maintenance debt.

Additionally, several crates were created for features that were absorbed into other crates,
never completed, or superseded by the norm system (\texttt{isls-norms}).

\textbf{This spec defines the clean target architecture and instructs Claude Code to
implement it in one pass---including removing all obsolete code.}

% ============================================================
\section{Crate Audit}
% ============================================================

Every crate is classified as \textcolor{keepgreen}{\textbf{KEEP}},
\textcolor{removered}{\textbf{REMOVE}}, or
\textcolor{modifyamber}{\textbf{MODIFY}}.

\subsection{KEEP --- Core Engine}

\begin{center}
\small
\begin{tabular}{>{\ttfamily}l l l}
\toprule
\textbf{Crate} & \textbf{Role} & \textbf{Action} \\
\midrule
isls-types & Core types: AppSpec, EntityDef, FieldDef, etc. & Keep as-is \\
isls-cli & CLI entry point, subcommands & Modify: remove obsolete cmds \\
isls-forge-llm & \textbf{THE engine}: HDAG, Staged Closure, structural, Oracle & Modify: apply \S\ref{sec:rules} \\
isls-hypercube & HyperCube, DomainRegistry, Fiedler/Kuramoto & Keep as-is \\
isls-norms & NormRegistry, 24 norms, composition, self-learning & Keep as-is \\
isls-evidence & SHA-256 evidence chain, audit trail & Keep as-is \\
isls-deployer & Docker deployment helpers & Keep as-is \\
\bottomrule
\end{tabular}
\end{center}

\subsection{KEEP --- Infrastructure (verify compiles)}

\begin{center}
\small
\begin{tabular}{>{\ttfamily}l l l}
\toprule
\textbf{Crate} & \textbf{Role} & \textbf{Action} \\
\midrule
isls-gateway & HTTP API for Studio dashboard & Keep, verify \\
isls-studio & Web UI (single HTML) & Keep, verify \\
isls-chat & Chat intent recognition & Keep, verify \\
isls-agent & No-code operator & Keep, verify \\
isls-reader & Barbara --- code reader/parser & Keep IF depended upon \\
isls-code-topo & Topology analyzer & Keep IF depended upon \\
\bottomrule
\end{tabular}
\end{center}

\subsection{REMOVE --- Blinddarm (Superseded or Absorbed)}

\begin{center}
\small
\begin{tabular}{>{\ttfamily}l p{8.5cm}}
\toprule
\textbf{Crate} & \textbf{Reason for Removal} \\
\midrule
isls-orchestrator & Tera templates superseded by structural generators in \texttt{isls-forge-llm} \\
isls-renderloop & Multi-pass render loop superseded by PCR Staged Closure \\
isls-blueprint & Crystal system superseded by \texttt{isls-norms} \\
isls-templates & Template files superseded by structural generators \\
isls-oracle & Oracle trait absorbed into \texttt{isls-forge-llm} \\
isls-decomposer & Decomposition absorbed into HDAG pipeline. Merge needed code into \texttt{isls-hypercube}. \\
isls-compose & Unclear purpose, likely absorbed \\
isls-capsule & Unclear purpose \\
isls-learner & Absorbed into \texttt{isls-norms} self-learning \\
isls-foundry & Absorbed or unused \\
isls-integrator & Absorbed or unused \\
isls-swarm & Separate feature (distributed). Remove from workspace. \\
isls-langgen & Schmiede. Separate feature. Remove from workspace. \\
isls-harness & Test harness. Keep ONLY if used by remaining tests. \\
isls-navigator & Unclear purpose \\
\bottomrule
\end{tabular}
\end{center}

\subsection{Removal Process}

\begin{enumerate}
\item For each REMOVE crate: check if any KEEP crate depends on it.
\item If yes: extract the needed code into the depending crate, \textbf{then} remove.
\item If no: delete the crate directory and remove from \texttt{Cargo.toml} workspace members.
\item After all removals: \texttt{cargo build --workspace} must succeed with 0 errors.
\item \texttt{cargo test --workspace} must pass.
\end{enumerate}

\textbf{CRITICAL:} Do \emph{not} leave dead workspace members. Do \emph{not} leave orphan directories.
The workspace after cleanup should have \textbf{10--14 crates}, not 25.

% ============================================================
\section{Target Architecture}
% ============================================================

\subsection{The Pipeline}

\begin{lstlisting}[language={}]
User Input (TOML or Chat)
    |
    v
+-------------------------------------+
|  isls-norms: NormRegistry            |  Match input to norms
|  -> ActivatedNorms                   |  Compose into plan
+-----------------+-------------------+
                  |
                  v
+-------------------------------------+
|  isls-hypercube: Decomposition       |  Fiedler eigenvalues
|  -> AppSpec with entities, fields,   |  Kuramoto synchronization
|     relationships                    |  Domain template matching
+-----------------+-------------------+
                  |
                  v
+-------------------------------------+
|  isls-forge-llm: HDAG Pipeline       |  THE ENGINE
|                                      |
|  S0 Ingest: AppSpec -> internal      |
|  S1 Canon:  normalize names          |
|  S2 Expand: build HDAG graph         |
|     (nodes = files, edges = deps)    |
|  S4 Solve:  topological traversal    |
|     structural -> deterministic      |
|     llm -> Oracle call               |
|  S5 Gate:   cargo check              |
|  S6 Coagula: fix errors (max 3)      |
|  S7 Emit:   output to disk          |
+-----------------+-------------------+
                  |
                  v
+-------------------------------------+
|  Output: Complete Rust project       |
|  backend/ + frontend/ +             |
|  docker-compose.yml                  |
|  -> cargo build -> docker-compose    |
|  -> Browser: working app             |
+-------------------------------------+
\end{lstlisting}

This is the \textbf{only} code generation path. No templates. No render loop.
Everything goes through \texttt{StagedClosure}.

\subsection{File Classification}

\subsubsection{Structural Files (Zero LLM Tokens, Zero Errors by Construction)}

\begin{center}
\small
\begin{tabular}{l l p{7cm}}
\toprule
\textbf{Layer} & \textbf{Files} & \textbf{Content} \\
\midrule
0 --- Infra & \texttt{.env.example}, \texttt{.gitignore}, \texttt{Cargo.toml}, \texttt{Dockerfile}, \texttt{docker-compose.yml}, \texttt{nginx.conf} & Generated from AppSpec metadata \\
0 --- Modules & \texttt{main.rs}, all \texttt{mod.rs} files & Module declarations + \texttt{pub use ::*} re-exports \\
1 --- Foundation & \texttt{errors.rs}, \texttt{pagination.rs}, \texttt{auth.rs}, \texttt{config.rs}, \texttt{pool.rs} & Fixed infrastructure types \\
3 --- Models & \texttt{models/\{entity\}.rs} & Derived from EntityDef: struct + Create/Update payloads \\
5 --- Migration & \texttt{001\_initial.sql} & CREATE TABLE + seed INSERT \\
6 --- Services & \texttt{services/\{entity\}.rs} & Thin delegation wrappers around queries \\
8 --- Frontend & \texttt{index.html}, \texttt{style.css}, \texttt{client.js}, \texttt{pages/*.js} & CRUD UI per entity \\
9 --- Tests & \texttt{api\_tests.rs} & Integration test stubs \\
\bottomrule
\end{tabular}
\end{center}

\subsubsection{LLM Files (Need Oracle, Caught by S6 Coagula)}

\begin{center}
\small
\begin{tabular}{l l p{7cm}}
\toprule
\textbf{Layer} & \textbf{Files} & \textbf{Content} \\
\midrule
5 --- Queries & \texttt{database/\{entity\}\_queries.rs} & SQL queries (SELECT, INSERT, UPDATE, DELETE) \\
7 --- Routes & \texttt{api/\{entity\}.rs}, \texttt{api/auth\_routes.rs} & HTTP request handlers \\
\bottomrule
\end{tabular}
\end{center}

\textbf{Ratio:} ${\sim}40$ structural : ${\sim}15$ LLM for a 7-entity application.

% ============================================================
\section{The Twelve Rules}
\label{sec:rules}
% ============================================================

Each rule was learned by burning a debug cycle. They are non-negotiable.

\subsection{Rule 1: Deterministic Content Must Be Structural}

If the file content is derivable from \texttt{EntityDef}, it \textbf{must not} be
LLM-generated. This includes: models, services, \texttt{mod.rs} files, migrations,
frontend pages. The LLM generates only SQL queries and HTTP route handlers---the
only places where genuine creativity adds value.

\subsection{Rule 2: \texttt{mod.rs} Must Re-Export with \texttt{pub use ::*}}

Both \texttt{services/mod.rs} and \texttt{database/mod.rs} must contain:
\begin{lstlisting}[language=C]
pub mod product;
pub use product::*;
\end{lstlisting}
This makes \emph{both} import patterns work---direct function imports and module aliases.
Neither causes E0603.

\subsection{Rule 3: Function Naming Is Fixed Across All Layers}

\begin{lstlisting}[language=C]
get_{entity}(pool, id)       -> Result<Entity, AppError>
list_{entities}(pool, params) -> Result<PaginatedResponse<Entity>, AppError>
create_{entity}(pool, payload) -> Result<Entity, AppError>
update_{entity}(pool, id, payload) -> Result<Entity, AppError>
delete_{entity}(pool, id)    -> Result<(), AppError>
\end{lstlisting}

No \texttt{\_service} suffix. No \texttt{\_handler} suffix. Same names in queries,
services, \emph{and} routes. The \texttt{ProvidedSymbol} signatures must list these
exact names.

\subsection{Rule 4: \texttt{AuthUser} Is a Function Parameter}

\begin{lstlisting}[language=C]
pub async fn list_products(
    pool: web::Data<PgPool>,
    params: web::Query<PaginationParams>,
    _user: AuthUser,  // extracted automatically by actix FromRequest
) -> Result<impl Responder, AppError> {
\end{lstlisting}

\textbf{Never} \texttt{req.extensions().get::<AuthUser>()}. It does not compile.

\subsection{Rule 5: Query Functions Take Owned Payloads}

\begin{lstlisting}[language=C]
pub async fn create_product(pool: &PgPool, payload: CreateProductPayload) -> ...
// NOT: payload: &CreateProductPayload
\end{lstlisting}

\subsection{Rule 6: Query Files Must Not Use \texttt{actix\_web} Types}

No \texttt{web::Data}, no \texttt{HttpResponse}, no \texttt{Responder}. Queries use only:
\texttt{sqlx::PgPool}, \texttt{sqlx::Row}, \texttt{crate::errors::AppError},
\texttt{crate::pagination::*}, \texttt{crate::models::*}.

\subsection{Rule 7: \texttt{bcrypt} Not \texttt{argon2}. \texttt{tracing} Not \texttt{log}.}

The \texttt{Cargo.toml} includes \texttt{bcrypt} and \texttt{tracing}. The LLM prompt
must explicitly state this.

\subsection{Rule 8: Seed INSERT Must Match Entity Fields}

The \texttt{generate\_migration()} function derives the seed INSERT from the User
entity's actual fields. No hardcoded \texttt{name} column.

\subsection{Rule 9: Dockerfile Uses \texttt{rust:1.88-slim}}

Current stable Rust is 1.88. Use this exact version for reproducibility.

\subsection{Rule 10: S6 Coagula Must Never Crash}

If error parsing fails, degrade gracefully---put all errors under an ``unknown'' key.
The Windows error format uses backslashes in paths. The parser must handle both
\texttt{/} and \texttt{\textbackslash}.

\subsection{Rule 11: \texttt{ProvidedSymbol} = Structural File Content}

Because the structural generator and the \texttt{ProvidedSymbol} factory use the
\emph{same} \texttt{EntityDef}, they produce identical types \textbf{by construction}.
No mismatch possible. This eliminated 90\% of errors.

\subsection{Rule 12: LLM Prompts Must Include Full Context}

Every LLM prompt must contain:
\begin{itemize}[nosep]
\item All available crate imports (\texttt{tracing}, \texttt{sqlx}, \texttt{serde}, \texttt{chrono}, \texttt{bcrypt})
\item The \emph{exact} model struct definitions (from structural generator)
\item The \emph{exact} function names to implement
\item Role-specific rules (QUERY: no actix types; ROUTE: AuthUser as parameter)
\item The ``REQUIRED FUNCTIONS'' list with exact signatures
\end{itemize}

% ============================================================
\section{Structural Generators}
% ============================================================

All generators reside in \texttt{isls-forge-llm/src/structural.rs}.
All take \texttt{\&AppSpec} or \texttt{\&EntityDef} and return \texttt{String}.
They are pure functions. No IO. No LLM. Deterministic.

\subsection{Required Generators}

\begin{lstlisting}[language=C]
// Layer 0 --- Infrastructure
fn generate_env_example(spec: &AppSpec) -> String;
fn generate_gitignore() -> String;
fn generate_cargo_toml(spec: &AppSpec) -> String;
fn generate_dockerfile() -> String;       // rust:1.88-slim
fn generate_docker_compose(spec: &AppSpec) -> String;
fn generate_main_rs(spec: &AppSpec) -> String;

// Layer 0 --- Module declarations (with pub use re-exports)
fn generate_models_mod(spec: &AppSpec) -> String;
fn generate_database_mod(spec: &AppSpec) -> String;
fn generate_services_mod(spec: &AppSpec) -> String;
fn generate_api_mod(spec: &AppSpec) -> String;

// Layer 1 --- Foundational types
fn generate_errors_rs() -> String;
fn generate_pagination_rs() -> String;
fn generate_auth_rs() -> String;
fn generate_config_rs() -> String;
fn generate_pool_rs() -> String;

// Layer 3 --- Models (per entity)
fn generate_model_rs(entity: &EntityDef) -> String;

// Layer 5 --- Migration
fn generate_migration(spec: &AppSpec) -> String;

// Layer 6 --- Services (per entity, thin delegation wrappers)
fn generate_service_rs(entity: &EntityDef) -> String;
fn generate_user_service_rs(entity: &EntityDef) -> String;  // +get_by_email

// Layer 8 --- Frontend
fn generate_frontend_index(spec: &AppSpec) -> String;
fn generate_frontend_style() -> String;
fn generate_frontend_client(spec: &AppSpec) -> String;
fn generate_frontend_page(entity: &EntityDef) -> String;
fn generate_nginx_conf() -> String;

// Layer 9 --- Tests
fn generate_api_tests(spec: &AppSpec) -> String;
\end{lstlisting}

\subsection{Model Generation Details}

\begin{lstlisting}[language=C]
fn generate_model_rs(entity: &EntityDef) -> String {
    // Main struct: #[derive(Debug, Serialize, Deserialize, FromRow, Clone)]
    //   Fields: id (i64) + entity fields (mapped types) + FK fields
    //           + created_at (DateTime<Utc>) + updated_at (DateTime<Utc>)
    //
    // CreatePayload: #[derive(Debug, Deserialize, Clone)]
    //   Fields: entity fields (respecting nullability) + FK fields
    //   NO id, NO created_at, NO updated_at (SystemGenerated)
    //
    // UpdatePayload: #[derive(Debug, Deserialize, Clone)]
    //   ALL fields as Option<T> (for partial updates)
}

fn field_to_rust_type(field: &FieldDef) -> String {
    // "String"|"string"|"text"|"varchar"  ->  "String"
    // "i32"|"int"|"integer"               ->  "i32"
    // "i64"|"bigint"                      ->  "i64"
    // "f64"|"float"|"double"              ->  "f64"
    // "bool"|"boolean"                    ->  "bool"
    // Wrap in Option<T> if field.nullable
}
\end{lstlisting}

\subsection{Service Generation Details}

\begin{lstlisting}[language=C]
fn generate_service_rs(entity: &EntityDef) -> String {
    // Imports: sqlx::PgPool, errors, models, pagination, queries
    //
    // 5 thin wrapper functions:
    //   pub async fn get_{snake}(pool: &PgPool, id: i64)
    //       -> Result<Entity, AppError> {
    //       {snake}_queries::get_{snake}(pool, id).await
    //   }
    //   ... list_{plural}, create_{snake}, update_{snake}, delete_{snake}
    //
    // Names: NO _service suffix. NO _handler suffix.
}
\end{lstlisting}

% ============================================================
\section{HDAG Prompt Templates}
% ============================================================

\subsection{Query File Prompt}

\begin{lstlisting}[language={}]
You are generating a Rust database query module.

## CRATE IMPORTS (include what you need):
use sqlx::PgPool;
use sqlx::Row;
use serde::{Serialize, Deserialize};
use chrono::{DateTime, Utc};
use tracing::{info, warn, error};
use bcrypt::{hash, verify};

## DO NOT USE:
- actix_web types (no web::Data, no HttpResponse, no Responder)
- log crate (use tracing instead)
- argon2 (use bcrypt instead)

## AVAILABLE TYPES:
{exact model struct from structural generator}
{exact CreatePayload struct}
{exact UpdatePayload struct}
{AppError enum definition}
{PaginationParams definition}
{PaginatedResponse definition}

## REQUIRED FUNCTIONS (use these EXACT names):
pub async fn get_{entity}(pool: &PgPool, id: i64)
    -> Result<{Entity}, AppError>
pub async fn list_{entities}(pool: &PgPool, params: &PaginationParams)
    -> Result<PaginatedResponse<{Entity}>, AppError>
pub async fn create_{entity}(pool: &PgPool, payload: Create{Entity}Payload)
    -> Result<{Entity}, AppError>
pub async fn update_{entity}(pool: &PgPool, id: i64,
    payload: Update{Entity}Payload) -> Result<{Entity}, AppError>
pub async fn delete_{entity}(pool: &PgPool, id: i64)
    -> Result<(), AppError>

For User entity, also:
pub async fn get_user_by_email(pool: &PgPool, email: &str)
    -> Result<User, AppError>

## QUERY RULES:
- Functions take OWNED payloads (not &CreatePayload)
- Use sqlx::query_as!() or sqlx::query() with .fetch_one/.fetch_all
- Map sqlx::Error to AppError::InternalError(e.to_string())
- For list: LIMIT/OFFSET from PaginationParams.offset() and .per_page()
- COUNT query for total in PaginatedResponse
- All functions are pub async fn

Return ONLY the Rust code. No markdown fences. No explanation.
\end{lstlisting}

\subsection{Route File Prompt}

\begin{lstlisting}[language={}]
You are generating Actix-web API route handlers.

## CRATE IMPORTS:
use actix_web::{web, HttpResponse, Responder};
use sqlx::PgPool;
use crate::auth::AuthUser;
use crate::errors::AppError;
use crate::pagination::PaginationParams;
use crate::models::{entity}::{{Entity}, Create{Entity}Payload,
    Update{Entity}Payload};
use crate::services::{entity} as {entity}_service;

## AUTH EXTRACTION:
AuthUser implements FromRequest. Use as FUNCTION PARAMETER:
  pub async fn list(..., _user: AuthUser) -> ...
NEVER use req.extensions().get::<AuthUser>()

## SERVICE CALLS:
  {entity}_service::get_{entity}(&pool, id).await
  {entity}_service::list_{entities}(&pool, &params).await
  {entity}_service::create_{entity}(&pool, payload.into_inner()).await

## RESPONSE PATTERN:
  Ok(HttpResponse::Ok().json(result))
  Ok(HttpResponse::Created().json(result))
  Ok(HttpResponse::NoContent().finish())  // for delete

Return ONLY the Rust code. No markdown fences. No explanation.
\end{lstlisting}

% ============================================================
\section{Verification Sequence}
% ============================================================

After implementation, Claude Code must execute this exact sequence:

\begin{lstlisting}[language=bash]
# 1. Workspace builds
cargo build --workspace --release
# Expected: 0 errors

# 2. All tests pass
cargo test --workspace
# Expected: 0 failures

# 3. Generate warehouse (mock)
rm -rf output/warehouse
cargo run --release -p isls-cli -- forge-v2 \
    --requirements examples/warehouse.toml \
    --mock-oracle --output ./output/warehouse

# 4. Generated code compiles
cd output/warehouse/backend && cargo build && cd ../../..
# Expected: 0 errors, <10 warnings

# 5. Generate with LLM (if API key available)
rm -rf output/warehouse_llm
cargo run --release -p isls-cli -- forge-v2 \
    --requirements examples/warehouse.toml \
    --api-key $OPENAI_API_KEY --output ./output/warehouse_llm

# 6. LLM-generated code compiles
cd output/warehouse_llm/backend && cargo build && cd ../../..
# Expected: 0 errors

# 7. Docker test
cd output/warehouse_llm
docker-compose up -d --build
sleep 30
curl -v http://localhost:8080/api/health
# Expected: 200 OK with JSON
docker-compose down
cd ../..

# 8. Verify workspace is clean
ls crates/ | wc -l
# Expected: 10-14 crates, not 25
\end{lstlisting}

% ============================================================
\section{Implementation Constraints for Claude Code}
% ============================================================

\begin{enumerate}
\item \textbf{Remove ALL obsolete crates first.} Before touching any code. Clean the workspace.
\item \textbf{If a removed crate is depended upon,} extract the needed code into the depending crate first.
\item \textbf{The \texttt{forge-v2} CLI command must still work.} Both \texttt{--mock-oracle} and \texttt{--api-key} modes.
\item \textbf{All structural generators go in} \texttt{isls-forge-llm/src/structural.rs}.
\item \textbf{All HDAG prompt builders go in} \texttt{isls-forge-llm/src/staged\_closure.rs} or a dedicated \texttt{prompts.rs}.
\item \textbf{All ProvidedSymbol factories go in} \texttt{isls-forge-llm/src/provided.rs}.
\item \textbf{No \texttt{unwrap()} in library code.}
\item \textbf{Every public item documented.}
\item \textbf{MIT License. Copyright (c) 2026 Sebastian Klemm.}
\item \textbf{Default LLM model: \texttt{gpt-4o}} (not \texttt{gpt-4o-mini}).
\item \textbf{Dockerfile: \texttt{rust:1.88-slim}.}
\item \textbf{Mock-oracle must produce compilable output.} This is the baseline. If mock works, LLM can only be equal or better.
\item \textbf{Commit after cleanup, commit after implementation, commit after verification.}
\end{enumerate}

% ============================================================
\section{Out of Scope (Future Work)}
% ============================================================

This spec is D1. The following are D2--D6 and are \emph{not} implemented here:

\begin{itemize}[nosep]
\item Norm-driven file selection (currently hardcoded for CRUD+Auth+Docker pattern)
\item Chat-to-App pipeline (\texttt{isls-chat} $\to$ \texttt{isls-agent} $\to$ forge)
\item Studio dashboard (\texttt{isls-studio} served by \texttt{isls-gateway})
\item Multi-domain proof (\texttt{ecommerce.toml}, \texttt{project.toml})
\item Incremental regeneration (amend plan, regenerate affected files only)
\item Claude API as Oracle alternative to OpenAI
\item Self-defining norm promotion (crystal $\to$ norm)
\item PSP-ASIC design for hardware-accelerated graph resonance
\end{itemize}

% ============================================================
\vfill
\noindent\rule{\textwidth}{0.4pt}
\begin{center}
\small
ISLS Final Architecture --- Sebastian Klemm --- March 2026\\
One spec. One pass. Clean workspace. Working pipeline.\\
HDAG sequences. Structural holds. LLM fills. Compiler verifies. Docker deploys.\\
\url{https://github.com/LashSesh/graphity}
\end{center}

\end{document}
